\chapter{The Ultimate Question}

For a long time, I believed the ultimate questions of life were these:

\begin{quote}
Who am I?

Where did I come from?

Where am I going?
\end{quote}

They sound deep, and they are not useless. Yet there is a problem hidden inside the first question. The self is not fixed. The person you become depends on what you do today, what you did yesterday, what you repeatedly pay attention to, and what you are willing to revise when reality proves you wrong.

So ``Who am I?'' can become a trap if it is treated as a final description. A better question is needed, one that can guide action while the self is still being built.

That question is:

\begin{quote}
What is more important?
\end{quote}

Ask it once, and it helps with a decision. Ask it repeatedly, honestly, and under pressure, and it begins to shape a value system. Ask it deeply enough, and it eventually produces the stricter question:

\begin{quote}
What is most important?
\end{quote}

This is one of the simplest tools in the whole book. It is also one of the most powerful. Wealth freedom is not reached by collecting clever ideas. It is reached by ranking things correctly, then acting according to the ranking long enough for compounding to appear.

\section*{Values Are Rankings}

A value system is not a list of beautiful words. It is an ordered list.

Many things are important. Money is important. Time is more important. Attention is more important still. Effort matters, but hard demand matters more. Persistence matters, but doing work with major significance matters more. Slides may matter in a presentation, but content matters more, and the person speaking may matter more than both.

Every serious choice contains an inequality.

\[
\text{attention} > \text{time} > \text{money}
\]

This small symbol, \(>\), is not decoration. It is a decision machine. When two attractive options compete, ask which belongs on the left side. When two fears compete, ask which one deserves more respect. When two duties compete, ask which one will still matter after time has passed.

Some chapters in this book were about equality: what a concept really means. What is capital? What is safety? What is knowledge? What is long-term? What is investment?

Other chapters were about multiplication: adding dimensions to personal competition, combining abilities, building growth rates, stacking advantages.

But this chapter is about inequality. It is about ranking.

\begin{center}
\begin{adjustbox}{max width=\linewidth}
\begin{tabular}{p{0.24\linewidth}p{0.58\linewidth}}
\hline
Symbol & Practical meaning \\
\hline
\(=\) & Define the concept correctly. \\
\(>\) & Rank what matters more. \\
\(\times\) & Combine dimensions so value compounds. \\
\hline
\end{tabular}
\end{adjustbox}
\end{center}

Most people do not fail because they have no values. They fail because their values are vague, unranked, or borrowed from the crowd. When the moment of choice arrives, the vague system collapses into impulse.

The question ``What is more important?'' sharpens the system before the pressure comes.

\section*{The Question Must Be Asked Again}

Because the self changes, the answer cannot be permanent.

A sincere answer today may become obsolete later. Not because it was dishonest, but because your experience, ability, obligations, and view of the future have changed. Growth often feels like being corrected by your own past. You discover that what once seemed important was only temporarily useful, or that what once seemed secondary was actually decisive.

This is not shameful. It is the normal cost of upgrading an operating system.

The danger is self-deception. A person may keep an old answer because it is familiar, because it protects pride, or because changing it would require action. That is how people become loyal to outdated versions of themselves.

So the method requires honesty:

\begin{enumerate}
\item Ask what matters.
\item Compare it with what matters more.
\item Keep comparing until the most important factor is exposed.
\item Act according to the answer.
\item Return later and ask again.
\end{enumerate}

A person who repeats this process becomes difficult to confuse. He may still make mistakes, but his mistakes become material for the next ranking.

\section*{Learning: The TOEFL Example}

When I prepared for TOEFL, my purpose was practical. I wanted to become a teacher at New Oriental because I needed a relatively stable and high income while my father was in the hospital.

At first, the obvious answer seemed to be vocabulary. Vocabulary is important. But the useful question was not ``Is vocabulary important?'' The useful question was:

\begin{quote}
Is anything more important than vocabulary?
\end{quote}

There had to be. Standardized tests can always include words that ordinary test-takers do not know. A rare technical word can defeat even a native speaker outside his field. If the whole strategy depends on knowing every word, then the strategy is fragile.

So the ranking changed.

The exam was not mainly testing the possession of an enormous word list. It was testing the ability to use basic vocabulary and context to infer meaning and logic. That realization changed the work. Instead of spending all my effort memorizing huge lists, I studied contextual reasoning inside reading passages.

That ranking later became a product: \emph{TOEFL Core Vocabulary: A 21-Day Breakthrough}. The book sold for years and became one of my early milestones on the road to wealth freedom.

The lesson is not that vocabulary is unimportant. The lesson is that the obvious important thing may not be the most important thing.

\[
\text{contextual reasoning} > \text{raw vocabulary size}
\]

One correct inequality can save years of waste.

\section*{Life: The Person You Can Reason With}

The same tool works outside study.

When choosing a life partner, many factors matter. Appearance matters. Education matters. Family background may matter. Temperament, habits, taste, health, and ambition may all matter. Pretending that they do not matter is false.

But again the useful question is stricter:

\begin{quote}
What matters more than these?
\end{quote}

After repeated comparison, I found the factor I considered most important:

\begin{quote}
Can this person reason?
\end{quote}

A person who can reason can discuss problems, correct mistakes, update views, negotiate differences, and face reality without turning every disagreement into a war. Once this condition is present, many other problems become solvable. If it is absent, even small problems can become permanent.

This is why the ranking matters.

\[
\text{capacity to reason} > \text{surface advantages}
\]

A marriage, partnership, or long collaboration is not protected by romance alone. It is protected by the ability to keep thinking together when thinking becomes uncomfortable.

\section*{Work: The Current Most Important Thing}

The same tool can govern a team.

A useful meeting can begin with one question:

\begin{quote}
What is the most important thing right now, and why?
\end{quote}

If the answer is clear, attention has a destination. If the answer is unclear, the meeting should not pretend to be productive. People may speak, report, argue, and plan, but without ranking they are only moving words around.

Management is often made mysterious. Many people imagine that leading a team requires a large collection of techniques. Some techniques are useful. But one ranking has helped me more than many techniques:

\begin{quote}
Choose a rapidly developing thing to work on.
\end{quote}

When a team is working on something that is growing quickly, many imperfections become tolerable. Everyone is busy learning, building, solving, serving, and responding to real demand. The work itself supplies energy and feedback.

When the thing grows slowly, problems multiply. Idle attention turns inward. Minor conflicts become major. Weak incentives become visible. Every defect appears larger because the work is not pulling people forward.

So before leading a team, the hardest work is often choosing the right thing for the team to do. If I cannot find the fast-developing direction, I would rather not begin. If I can find it, many later management problems become smaller.

\[
\text{rapidly developing direction} > \text{management technique}
\]

\section*{Investing: The Most Important Action May Be Inaction}

Investment is where this question becomes severe.

After entering the investment world, I had to ask again:

\begin{quote}
What is more important? In the end, what is most important?
\end{quote}

The answer I reached was plain:

\begin{quote}
After buying an appreciating asset with sustainable long-term growth, hold it. Do not move unnecessarily.
\end{quote}

This sounds simple only after the hard thinking is finished.

In the mirror world of investing, much of the real effort must happen before money is committed. One must understand the asset, the people behind it, the growth rate, the risk, the position size, the time horizon, and one's own emotional limits. After that, frequent action often becomes damage.

This ranking produces several principles:

\begin{itemize}
\item Be responsible for your own decisions.
\item Growth rate determines value growth.
\item Invest in people stronger than yourself.
\item Finish the homework before the money leaves your hand.
\item Do not enter what you do not understand, no matter how attractive it looks.
\item In finance, capital itself is a force that must be respected.
\end{itemize}

The investor's inequality is:

\[
\text{right asset held long enough} > \text{clever movement}
\]

Doing nothing after a correct decision is not laziness. It is discipline. It is the ability to protect a good ranking from emotional noise.

\section*{Foresight}

People with foresight ask different questions.

While others write an article, they may be writing a book. While others live entirely in the present, they may be practicing life in the future. While others upgrade software, they may be upgrading the operating system.

This difference compounds. The person who asks higher-level questions begins from a different place. He does not merely work harder inside the old frame. He changes the frame.

For a long time, wealth freedom appeared to be the destination. Later, after crossing certain cost lines, it became clear that personal wealth freedom was only a beginning. Beyond it are family wealth accumulation, wealth management, and eventually inheritance.

Then the question appears again:

\begin{quote}
In inheritance, what is more important?
\end{quote}

The answer is not merely money. Money can be lost by people who lack the ability to manage it. Status can decay. Assets can be divided, wasted, or misunderstood.

The more important inheritance is ability.

\[
\text{ability} > \text{assets alone}
\]

The ability to learn, judge, cooperate, invest, create, control attention, and act responsibly is more durable than any single asset. A family that transfers money without transferring ability may only delay decline. A family that transfers ability has a chance to renew itself.

\section*{The Learner Is the Decisive Factor}

This also applies to education.

In learning, the teacher matters. The environment matters. Materials, peers, tools, and timing matter. But when the comparison is pushed far enough, the learner becomes the decisive factor.

The same teacher teaches many students. The same environment surrounds many people. Yet some rise faster than others. The difference cannot be explained only by external conditions.

So the ranking is:

\[
\text{self-directed learner} > \text{environment} > \text{teacher}
\]

This does not insult teachers or environments. It locates responsibility. A good teacher can open doors, but the student must walk. A good environment can reduce friction, but the student must use it. A good book can contain a tool, but the reader must install it into behavior.

A new tool unused is like a powerful phone treated only as a device for calls and messages. The hardware may be excellent, but the user receives little of its value. Concepts are the same. Reading about attention, capital, growth rate, hard demand, long-term holding, or personal business models changes nothing unless the ideas are used.

\section*{The Most Important Personal Quality}

If the learner is decisive, then another question follows:

\begin{quote}
What personal quality matters most?
\end{quote}

Courage matters. Patience matters. Intelligence matters. Toughness matters. But one quality has proven more important than most people expect:

\begin{quote}
Do the work, then learn to love it.
\end{quote}

Many people know how to love a thing first and then work on it. That is easy. The harder and more valuable ability is to take a necessary field, begin working, improve inside it, and gradually build affection for it.

Human beings rarely do excellent work in things they hate. But love is not always a precondition. Often it is a result.

People frequently say, ``I have no interest in my major,'' or ``I have no interest in this work.'' In many cases, the deeper truth is simpler: they have not yet become good enough to enjoy it. When progress is absent, the work feels hostile. When progress appears, interest often follows.

The rough method is:

\begin{enumerate}
\item Learn the minimum necessary knowledge.
\item Practice immediately.
\item Endure the clumsy beginning.
\item Notice real improvement.
\item Let improvement create interest.
\item Let interest support further practice.
\end{enumerate}

\begin{center}
\begin{adjustbox}{max width=\linewidth}
\begin{tikzpicture}[node distance=0.8cm, every node/.style={font=\small, align=center}]
\node[draw, rounded corners, minimum width=2.45cm, minimum height=0.72cm] (know) {minimum\\knowledge};
\node[draw, rounded corners, minimum width=2.45cm, minimum height=0.72cm, right=of know] (practice) {practice};
\node[draw, rounded corners, minimum width=2.45cm, minimum height=0.72cm, right=of practice] (progress) {visible\\progress};
\node[draw, rounded corners, minimum width=2.45cm, minimum height=0.72cm, right=of progress] (love) {interest\\and love};
\draw[->] (know) -- (practice);
\draw[->] (practice) -- (progress);
\draw[->] (progress) -- (love);
\draw[->] (love) to[bend left=35] (practice);
\end{tikzpicture}
\end{adjustbox}
\end{center}

This is why ``persistence'' and ``willpower'' can be overvalued. A person who gives work enough meaning and then improves inside it may not experience the process as grim persistence. He may experience it as a relationship.

I often treat skills this way. I do not merely use them. I cultivate them, return to them, repair them, and let them become part of my life. When a skill becomes beloved, outside resistance loses power.

Execution can then be judged by a simple test:

\begin{quote}
When the person is not yet doing well, does he continue?
\end{quote}

Anyone can continue when praise, money, and progress arrive quickly. Execution is revealed before those rewards arrive.

\section*{The Slide Deck Is Not the Point}

The same ranking works even in small practical matters.

Many people spend large amounts of time learning slide design. Good slides are useful. Beauty is not wrong. But in a presentation, are slides the most important thing?

No. Content is more important.

\[
\text{content} > \text{slide design}
\]

Reid Hoffman once released the slide deck used for LinkedIn's 2004 Series B pitch. The company raised \(\$10\) million from Greylock. By ordinary design standards, the slides were unimpressive. Yet the pitch worked.

Why?

Investors were not mainly buying fonts, colors, or decorative polish. They were evaluating facts, logic, market structure, network effects, and the person speaking. More importantly, the speaker was Reid Hoffman: Stanford graduate, Oxford philosophy master's holder, PayPal executive, experienced founder, and already a deeply connected person in Silicon Valley.

So the ranking must be improved:

\[
\text{person} > \text{content} > \text{slide design}
\]

This does not mean slide design is worthless. It means the highest-return work is often elsewhere.

If I need beautiful slides, I can buy a good template for 15--30 dollars and save hours. The more serious question is not ``How do I become slightly better at slide decoration?'' It is:

\begin{quote}
How do I become the person whose words deserve attention?
\end{quote}

This question is uncomfortable because it refuses cosmetic improvement. It points back to substance: track record, judgment, courage, insight, trust, and contribution.

A person with little weight may speak correctly and still be ignored. A person with real weight may speak plainly and still be heard. So the work is not only to prepare better sentences. The work is to become someone whose sentences carry evidence.

\section*{A Pocket Decision Method}

The question ``What is more important?'' can be carried into almost any decision.

\begin{center}
\begin{adjustbox}{max width=\linewidth}
\begin{tabular}{p{0.34\linewidth}p{0.48\linewidth}}
\hline
Situation & Better ranking question \\
\hline
Studying & What ability actually drives the result? \\
Partner choice & What quality makes future problems solvable? \\
Team work & What is the current most important thing? \\
Investing & What must be decided before action, and what should be left alone after? \\
Education & Which part depends on me? \\
Career & What makes me the person who can say this? \\
Inheritance & What ability must survive beyond the asset? \\
\hline
\end{tabular}
\end{adjustbox}
\end{center}

The method is not complicated:

\begin{enumerate}
\item List what seems important.
\item Admit that many items really are important.
\item Compare them in pairs.
\item Move the more important item to the left.
\item Continue until the highest-order factor appears.
\item Build action around that factor.
\end{enumerate}

This protects attention. Without ranking, every important thing competes equally, and attention is torn apart. With ranking, attention can concentrate.

The entire path to wealth freedom depends on concentrated attention. Capital compounds only when attention has first been allocated well enough to produce skill, judgment, assets, and patience.

\section*{The Tool Must Be Used}

A tool received is not the same as a tool mastered.

Concepts can be handed over through writing. Methods can be explained. Examples can be shown. But the repeated use of the tool belongs to the reader. No teacher can ask the question for you at the decisive moment. No book can force you to rank correctly when vanity, fear, desire, or the crowd pushes in another direction.

This is why responsibility eventually returns to the self.

A person can read about attention and still waste it. He can read about capital and still disrespect it. He can read about long-term holding and still trade from anxiety. He can read about hard demand and still build products for his own fantasy. He can read about the ultimate question and still avoid asking it when the answer would be inconvenient.

The question is a blade. It separates the merely important from the more important, and the more important from the most important.

Use it on study. Use it on work. Use it on relationships. Use it before investing. Use it after success. Use it when choosing what to ignore. Use it when deciding who to become.

The path does not become easy. It becomes clearer.
